\subsection{Base Conversion}
\label{sec:rns-base-conversion}
%
\begin{definition}[RNS Basis] 
    Given a large modulus $Q = \prod_{i=0}^{k-1} q_i$ composed of pairwise coprime factors, the \textit{RNS basis} is defined as $\mathcal{B}_Q = \{q_0, q_1, \dots, q_{k-1}\}$ as the set of RNS moduli used in the RNS/CRT representation \eqref{form:crt}.
    \begin{equation}
        \mathbb{Z}_Q \cong \mathbb{Z}_{q_0} \times \mathbb{Z}_{q_1} \times \cdots \times \mathbb{Z}_{q_{k-1}}
    \end{equation}
\end{definition}

% A subset $\mathcal{B}_{Q} \subseteq \mathcal{B}$ forms the basis for the $Q$-towers of...

%%
%% BASE EXTENSION (NAIVE)
\begin{definition}[Base conversion]
    Let $\mathcal{B}_Q = \{q_0,\dots,q_{k-1}\}$ and
    $\mathcal{B} = \{b_0,\dots,b_{\ell-1}\}$ be RNS bases with
    $\gcd(Q, B) = 1$, where $B = \prod_{j} b_j$. Base conversion is
    the map that takes the representation of $x$ in $\mathcal{B}_Q$
    to the representation of $[x]_Q$ in $\mathcal{B}$, i.e.
    \begin{equation*}
        \left([x]_{q_0}, \dots, [x]_{q_{k-1}}\right)
        \;\longmapsto\;
        \left([x]_{b_0}, \dots, [x]_{b_{\ell-1}}\right)
    \end{equation*}
    When the residues in $\mathcal{B}$ are appended to those in $\mathcal{B}_Q$, so that $x$ becomes represented in the \textit{extended} basis $\mathcal{B}_Q \cup \mathcal{B}$ for the modulus $QB$, the operation is also referred to as \emph{base extension}.
\end{definition}

% \begin{definition}[Base Extension.]
%     Special case of base conversion where we convert from $\mathcal{B}_Q \longmapsto \mathcal{B}_{QP}$. Note $\mathcal{B}_{QP} = \mathcal{B}_{Q} \cup \mathcal{B}_{P}$ and $P = \prod_i p_i$ etc. 
% \end{definition}

\begin{remark}
    \color{purple}
    Extend mod notation for RNS polynomials: $[a]_{\mathcal{B}_Q} = ([a]_{q_0}, \dots, [a]_{q_{k-1}})$ for the basis $\mathcal{B}_Q = \{q_i\}_{i=0}^{k-1}$ means $a$ must be in RNS form 
    % (\textcolor{red}{TODO: distinguish CRT vs NTT mode... maybe $[a]^\mathsf{NTT}_{\mathcal{B}_Q}$?})
\end{remark}

The naive approach to base conversion reconstructs $a$ from its residues through inverse \gls{crt}, and then reduces the result against each modulus of the new basis. This reintroduces the multi-precision arithmetic that RNS was meant to avoid. % Cheaper but approximate alternatives exists, described next.


\subsubsection{Fast Base Conversion}
Introduced by \cite{cryptoeprint:2016/510}, fast base conversion operates directly on the residues in CRT mode, avoiding the CRT reconstruction and multi-precision integers entirely. \textcolor{red}{(Does not work in evaluation mode).} The price is approximation, since \eqref{eq:fastbconv} produces not $x$ but $x + \alpha_x Q$ for some integer $\alpha_x \in [0, k)$. \textcolor{purple}{Computing the exact overflow $\alpha_x$ is possible but would require costly operations in RNS, and is unnecessary in \textit{extension} by $P$.} \textcolor{red}{Introduce the notation $\hat{q}_i = Q/q_i$ here?}

\begin{equation}
    \label{eq:fastbconv}
    \mathsf{FastBConv}(x, \mathcal{B}_Q, \mathcal{B}) = \left(\left[\sum_{i=0}^{k-1}\left[x_i\left(\frac{Q}{q_i}\right)^{-1}\right]_{q_i} \times \frac{Q}{q_i}\right]_b\right)_{b\in \mathcal{B}}
\end{equation}


\subsubsection{Fast Base Extension}
\textcolor{red}{We can use} \eqref{eq:fastbconv} as a subroutine in base extension. Only the new residues $\mathcal{B}_Q \longmapsto \mathcal{B}_P$ need to be calculated, and since $\mathcal{B}_{QP} = \mathcal{B}_Q\cup \mathcal{B}_P$ the original residues are kept unchanged.

%%%
%%% Fast Base Extension
%%%
\begin{algorithm}[ht!]
\caption{Fast Base Extension}
\label{alg:fast_base_extension}
\begin{algorithmic}[1]
    \Require $\NTT{[a]}_{\mathcal{B}_Q}$ and a basis $\mathcal{B}_P$
    \Ensure $\NTT{[a + \epsilon]}_{\mathcal{B}_{QP}}$ where $\epsilon = \alpha_aQ$ for some $\alpha_a\in[0, k)$
    % \State $\mathsf{shared} \gets 0$
    % \For{$c \in \Call{InverseNTT}{a}$}
    %     \Comment{Convert to CRT form}
    %     \State $\mathsf{shared} \gets \mathsf{shared} + [c\cdot\hat{q}_i^{-1}]$ 
    % \EndFor
    \State $a_{Q} := \Call{InverseNTT}{a}$
    \State $\mathsf{shared} \gets \sum_{i = 0}^{k} [a_Q[i]\hat{q}_i^{-1}]_{q_i}\cdot \hat{q}_i$
    \For{$p \in \mathcal{B}_P$}
        \State $a \gets a \cup \Call{NTT}{[\mathsf{shared}]_p}$
    \EndFor
    \State\Return $a$
\end{algorithmic}
\end{algorithm}

% \State $a_{Q} := \Call{InverseNTT}{a}$
% \State $a_{P} := \mathsf{FastBConv}(a_Q, \mathcal{B}_Q, \mathcal{B}_P)$
% \State\Return $a \cup \Call{NTT}{a_{P}}$ 

% Notably, \eqref{eq:fastbconv} is approximate as a \textit{conversion} but exact as an \textit{extension}: the extended residue vector over $\mathcal{B}_Q\cup\mathcal{B}_P$ is an exact RNS representation of the lifted integer $x' = [x]_Q + \alpha_x Q \in \mathbb{Z}_{QP}$, since $\alpha_x Q \equiv 0 \pmod{q_i}$ keeps the original (reused) $Q$-residues consistent with $x'$. The RNS schemes are designed to tolerate this overflow, which is removed --- or shrunk to the small $\alpha_x$ --- by a subsequent division during mod down.

\begin{remark}
    The factors $\hat{q}^{-1}_i =[(\frac{Q}{q_i})^{-1}]_{q_i}$ and $\hat{q}_i =\frac{Q}{q_i}$ can be pre-calculated and stored in memory (discussed in \autoref{sec:impl:core}).
\end{remark}

\subsubsection{Scale}
Note that the error $\alpha_xQ \pmod{QP}$ produced by \autoref{alg:fast_base_extension} would be completely removed if we multiplied by $P$. This observation leads to the specialized operation ``\textit{extend and scale by} $P$''. I refer to this operation as $\Call{Scale}{x, P} \mapsto [xP]_{\mathcal{B}_{QP}}$. 
\begin{gather*}
    x' = x + \alpha_xQ \pmod{QP} \\
    x'P = xP + \alpha_x QP = xP \pmod{QP}
\end{gather*}
\begin{equation}
    \label{def:scale}
    \Call{Scale}{x, P} = \Call{FastBaseExtend}{x, \mathcal{B}_Q, \mathcal{B}_P} \cdot P
\end{equation}

Furthermore, in \gls{rns}, the $P$-towers ($\bmod{q_i} \;\forall\; i$) all become 0 when multiplied by $P = p_0\dots p_{k'-1}$. This operation is therefore far cheaper in reality as it can be executed by only appending $k'$ empty towers with the $p_i$ moduli to $x$, and multiplying each $Q$-tower by $P$.

\begin{algorithm}[ht!]
\caption{\protect\Call{Scale}{x, P}}
\label{alg:scale}
\begin{algorithmic}[1]
    \Require $\NTT{[x]}_{\mathcal{B}_Q}$ with $\|\mathcal{B}_{Q}\| = k_q$ and a basis $\mathcal{B}_P$ with $\|\mathcal{B}_{P}\| = k_p$
    \Ensure $\NTT{[xP]}_{\mathcal{B}_{QP}}$
    \For{$(x_i, q_i)\in x$}
        \State $x_i \gets x_i \cdot P \bmod{q_i}$
    \EndFor
    \For{$i \in [0, k_p)$}
        \State $x_{k + i} \gets 0 \bmod{\mathcal{B}_P[i]}$
        \Comment{Append empty $p_i$-tower to $x$}
    \EndFor
    \Return $x$
\end{algorithmic}
\end{algorithm}

\subsubsection{\color{red} Modulus Switching}
As rounding is not natively supported in \gls{rns}, the textbook approach to modulus switching is replaced by dropping the last limb and applying a small correction to each limb. 

\begin{equation}
    [c']_{q_i} = [q_{k-1}^{-1}(c - \delta)]_{q_i}, \quad 0 \leq i < k - 1
\end{equation}



% Rounding is not natively supported in \gls{rns} because the independent, non-positional nature of the residue limbs obscures the overall magnitude, making fractional corrections impossible without full Chinese Remainder Theorem (CRT) reconstruction. Consequently, to emulate the modulus switch—scaling down from $Q$ to $Q^{\prime}=Q/q_{k-1}$—the classical approach must be adapted. In an \gls{rns} architecture, this fundamentally amounts to dropping the last limb $q_{k-1}$. However, to mathematically replicate the rounding step and preserve decryption correctness, a precise correction term $\delta\equiv c \pmod{q_{k-1}}$ must be applied to the remaining limbs. The exact \gls{rns} modulus switch is therefore computed limb-wise as:  

% When operating within a Residue Number System (RNS), the classical modulus switch—scaling down from $Q$ to $Q^{\prime}=Q/q_{k-1}$—amounts to dropping the last RNS limb $q_{k-1}$. However, to mathematically replicate the "rounding" step of Definition 2.7 and maintain correct decryption, a small correction term $\delta \equiv c \pmod{q_{k-1}}$ must be applied to the remaining limbs. The exact RNS modulus switch is computed limb-wise as:  $$[c^{\prime}]_{q_{i}}=[q_{k-1}^{-1}(c-\delta)]_{q_{i}}$$for $0\le i<k-1$. For the BGV scheme specifically, the correction must additionally satisfy $\tilde{\delta}\equiv0 \pmod{t}$ to ensure the plaintext message is maintained. This exact operation successfully shrinks the noise by a factor of $q_{k-1}$.



% Modulus switching is performed 

% In leveled homomorphic encryption, modulus switching is the primary tool for noise management, converting a ciphertext $\vec{x} \in R_{Q}^{2}$ to a smaller ciphertext modulus $Q^{\prime}$ while preserving the underlying plaintext after decryption. In a standard setting, this is achieved by scaling the message by a factor of $Q^{\prime}/Q$ and rounding to the nearest integer. 

% % This process proportionally scales down the intrinsic noise, permitting further homomorphic computation without the need for bootstrapping

% In RNS, modulus switching --- also referred to as \textit{mod down} --- from $Q$ to $Q' = Q/q_{k-1}$ amounts to \textit{dropping the last limb} with a small correction, computed limb-wise without CRT reconstruction:
% \begin{equation}
%     \label{eq:rns-modswitch}
%     [c']_{q_i} = \left[q_{k-1}^{-1}\left(c - \delta\right)\right]_{q_i}, \quad 0 \leq i < k-1,
% \end{equation}
% where $\delta \equiv c \pmod{q_{k-1}}$ and, for BGV, $\delta \equiv 0 \pmod{t}$ so that decryption correctness is maintained. The noise shrinks by a factor $q_{k-1}$, matching the classical modulus switch.

\subsubsection{\color{red} Approximate Mod Down}
Modulus switching exactly requires multiprecision integers.. can use FBE again... exploits the fact that the polynomial is pre-multiplied by $P$...

Uses $\delta = \Call{FastBConv}{[x]_{\mathcal{B}_P}, \mathcal{B}_P, \mathcal{B}_Q}$
% The reverse operation, switching from the extended modulus $QP$ back down to $Q$, uses \eqref{eq:fastbconv} to convert the $P$-part back to $\mathcal{B}_Q$ before dividing by $P$ \autocite{cryptoeprint:2012/099, cryptoeprint:2018/117}:
\begin{equation}
    \label{eq:approxmoddown}
    \Call{ApproxModDown}{x} = \left[P^{-1}\left(x - \Call{FastBConv}{[x]_{\mathcal{B}_P}, \mathcal{B}_P, \mathcal{B}_Q}\right)\right]_{\mathcal{B}_Q}
\end{equation}
For BGV the converted term must additionally be corrected modulo $t$ to preserve the plaintext \cite{cryptoeprint:2021/204}. % This operation also corrects the fast base conversion error: the conversion overflow $\alpha P$ is divided by $P$, shrinking it to the small additive term $\alpha$ --- which is why the approximation in \eqref{eq:fastbconv} is tolerable.
